\section{Introduction}
\label{sec:intro}

The spatial distribution of stellar ages, metallicities, and
$\alpha$-element abundances within early-type galaxies (ETGs) encodes
the history of their assembly and quenching
\citep{Thomas2005,Kuntschner2010}.
Negative metallicity gradients, long interpreted as a hallmark of
monolithic dissipative collapse \citep{Larson1974}, persist even in
massive systems that have experienced significant merger activity
\citep{Greene2015}.
More recently, integral-field unit (IFU) surveys---most notably
ATLAS$^{\mathrm{3D}}$, CALIFA, MaNGA, and SAMI---have opened the
door to systematic, two-dimensional mapping of stellar populations
across representative galaxy samples \citep{McDermid2015,GonzalezDelgado2015}.

A consistent picture has nonetheless proved elusive.
The slope of the age gradient depends sensitively on both the
stellar-population model and the assumed star-formation history (SFH)
parametrisation \citep{Goddard2017}.
Furthermore, whether massive ETGs built up their mass
\emph{inside-out} (central regions forming first) or
\emph{outside-in} (merger-driven outer-shell growth) remains
debated \citep{Peletier1990,Trujillo2011}.

In this paper we revisit stellar-population gradients using
32 ETGs from the SAMI Galaxy Survey and address three specific
questions:
(i) How steep are the metallicity and age gradients as a function of
    stellar mass?
(ii) Does the $\afe$ profile vary with radius, providing clues to
     the spatial variation of star-formation time-scale?
(iii) Are the gradients consistent with an inside-out quenching
     scenario?

The paper is organised as follows.
Section~\ref{sec:data} describes the sample and data reduction.
Section~\ref{sec:methods} outlines the full-spectrum fitting
procedure.
Section~\ref{sec:results} presents the measured gradients.
Section~\ref{sec:discussion} discusses implications for galaxy
formation, and we conclude in Section~\ref{sec:conclusions}.
Throughout we assume a flat $\Lambda$CDM cosmology with
$H_0 = 70\,\mathrm{km\,s^{-1}\,Mpc^{-1}}$, $\Omega_\mathrm{m}=0.3$.

%%====================================================================
